\begin{figure}[h]
    \centering
    \begin{subfigure}[b]{.49\textwidth}
        \includegraphics[width=0.99\textwidth]{figs/cars_confidence_boxplot.png}
        \caption{Cars dataset}
    \end{subfigure}
    \begin{subfigure}[b]{.49\textwidth}
        \includegraphics[width=0.99\textwidth]{figs/pets_confidence_boxplot.png}
        \caption{Pets dataset}
    \end{subfigure}
    \caption{Boxplot of confidence values when prompts with varying specificity are used. \textbf{CLS} refers to using class names from the dataset for prompting. \textbf{ZS} are prompts specific to each dataset described in more detail in Appendix \ref{sec:prompt_specificity}. \textbf{GPT3}, \textbf{Wikidef}, \textbf{Hi} and \textbf{Worddef} are prompt templates taken from \cite{li2022elevater}. \textbf{GPT3} refers to descriptions of each class generated by OpenAI-GPT3 \cite{pretrain-zeroshot} and used as the prompt. \textbf{Worddef} and \textbf{Wikidef} are Wordnet \cite{wordnet} and Wiktionary \cite{wiktionary} definitions of each class in the dataset (With the exception of Cars for which no classes have Wordnet definitions). \textbf{Hi} refers to a set of words extracted from Wordnet based on the ontological hierarchy of the class name. } %\CRO{1. use only one legend. 2. increase font size. 3. remove the value of $\tau_R$ and change it for loose and strict refinement 4. Better to be prepared to use just one of those plots (Cars maybe) we need to add the plot of $\tau_R$ vs accuracy}} %\CRO{We have to do these figures again}}
    \label{fig:confidence_boxplot}
    % \vspace{-4.5em}
\end{figure}